\section{总结与延伸}

\begin{enumerate}
    \item 奖励设计是 RL 应用的瓶颈，reward learning 提供了自动化的解决方案。
    \item \textbf{Goal Classifiers}：从成功/失败示例学习，简单直接。对抗训练防止 reward hacking。
    \item \textbf{Inverse RL}：从专家行为推断奖励函数，比模仿学习更灵活。
    \item \textbf{偏好学习}：从人类偏好比较中训练奖励模型，是 RLHF 的核心。
    \item 所有 reward learning 方法都面临 reward hacking 的风险，需要正则化和约束。
\end{enumerate}

\subsection{方法对照表}

\begin{table}[H]
\centering
\begin{tabular}{p{0.18\textwidth} p{0.25\textwidth} p{0.25\textwidth} p{0.22\textwidth}}
\toprule
方法 & 核心信号 & 主要优势 & 主要风险 \\
\midrule
Goal Classifier & 成功 vs 失败状态的二分类概率 & 简单、易于小数据集应用 & Reward hacking、分布漂移 \\
Inverse RL & 专家轨迹优化的 reward 函数 & 可以解释专家意图 & 需要高质量轨迹、训练不稳定 \\
偏好学习 (RLHF) & 人类选择的 winner vs loser & 直接对齐人类喜好 & Reward hacking、评分噪声 \\
\bottomrule
\end{tabular}
\caption{Reward learning 方法对照}
\end{table}

\subsection{延伸阅读与实践方向}

\begin{itemize}
    \item 继续跟踪 Stanford RL group 的 reward modeling workshop 记录
    \item 学习《Deep RL from Human Preferences》中的数据收集 protocol
    \item 把自己的 logging pipeline 和 reward trace 绑定在一起，建立可回滚的版本 control
\end{itemize}

\subsection{拓展阅读}
\begin{itemize}
    \item Christiano et al., \textit{Deep Reinforcement Learning from Human Preferences} (2017)。
    \item Ziegler et al., \textit{Fine-Tuning Language Models from Human Preferences} (2019)。
    \item Sharma et al., \textit{Self-Improving Robots: End-to-End Autonomous Visuomotor Reinforcement Learning} (2023)。
\end{itemize}

\end{document}
